\documentclass[12pt]{article}
\usepackage[T1]{fontenc}
\usepackage[utf8]{inputenc}
\usepackage{lmodern}
\usepackage{textcomp}
\usepackage{enumitem}
\usepackage{geometry}
\usepackage{graphicx}
\usepackage{amsmath, amssymb}
\geometry{margin=1in}
\begin{document}
\begin{center}
Chemistry - Undergraduate Level Assessment 24 \
Total Marks: 40 \
Answer all questions.
\end{center}

\section*{Multiple Choice (10 marks)}
\begin{enumerate}[label=\arabic*.]
\item Which of the following is lower for argon than for neon?
\begin{enumerate}[label=(\alph*)]
    \item Melting point
    \item Boiling point
    \item Polarizability
    \item First ionization energy
\end{enumerate}
\item When the following equation is balanced, which of the following is true? \_\_ MnO 4- + \_\_ I- + \_\_ H+ \textless{}-\textgreater{} \_\_ Mn$_{2}$+ + \_\_ IO 3- + \_\_ H$_{2}$O
\begin{enumerate}[label=(\alph*)]
    \item The I- : IO 3- ratio is 3:1.
    \item The MnO 4- : I- ratio is 6:5.
    \item The MnO 4- : Mn$_{2}$+ ratio is 3:1.
    \item The H+ : I- ratio is 2:1.
\end{enumerate}
\item Calculate the Larmor frequency for a proton in a magnetic field of 1 T.
\begin{enumerate}[label=(\alph*)]
    \item 23.56 GHz
    \item 42.58 MHz
    \item 74.34 kHz
    \item 13.93 MHz
\end{enumerate}
\item Of the following compounds, which has the lowest melting point?
\begin{enumerate}[label=(\alph*)]
    \item HCl
    \item AgCl
    \item CaCl 2
    \item CCl 4
\end{enumerate}
\item Of the following compounds, which is LEAST likely to behave as a Lewis acid?
\begin{enumerate}[label=(\alph*)]
    \item BeCl 2
    \item MgCl 2
    \item ZnCl 2
    \item SCl 2
\end{enumerate}
\end{enumerate}

\section*{True / False (10 marks)}
\textit{Answer True or False}
\begin{enumerate}[label=\arabic*.]
\setcounter{enumi}{5}
\item True or False: Tetramethylsilane is commonly used as a 1 H reference because its protons are weakly shielded in NMR spectra.
\item True or False: The hybrid $sp^3$ orbitals for carbon are formed from three atomic orbitals.
\item True or False: The 13 C NMR spectrum of 2-methylpentane exhibits five distinct chemical shifts due to its unique structural isomers.
\item True or False: Protons in a molecule with a rotational correlation time of 1 ns have the fastest spin-lattice relaxation at a magnetic field of 3.74 T.
\item True or False: A$_{0}$.100 M HCl solution has a higher ionic strength than a 0.050 M CaCl 2 solution because it is a strong acid.
\end{enumerate}

\section*{Long Form Response (20 marks)}
\textit{Answer each question with a clear and complete explanation.}
\begin{enumerate}[label=\arabic*.]
\setcounter{enumi}{10}
\item Discuss the factors that contribute to the electron affinity of elements and analyze why calcium exhibits the lowest electron affinity among its peers.
\item Discuss the concept of conjugate acid-base pairs, providing an example, and analyze the significance of the pair H$_{3}$O+ and H$_{2}$O in acid-base chemistry.
\end{enumerate}

\vfill
\noindent\textit{End of Paper}
\end{document}